%! Regression Lines and Residuals — Unit 2, Lesson 2.3 — Homework (Answer Key)
% GRADUAL RELEASE: this is the lesson's SCORED practice and the LAST component in the packet.
% Paper homework is the DEFAULT for this course; DeltaMath is an occasional override that
% replaces this page and strikes its score cell on the cover. Started in the last eight minutes
% of the period and finished after two study halls.
% THIRD CONTEXT: the notes teach on pizza deliveries, Guided Practice runs on sleep and reaction
% time, and these items are on used-car prices — recall is not the skill being built.
% ONE PAGE: five items in two parts (A1–A3, B1–B2), then the spiralbox preview of Lesson 2.4.
% A2 and B2 carry the sign of the residual — the lesson's target misconception.
% \answerspace{H}{} reserves exactly H in the blank; the key passes the answer as the second
% argument, so the two files paginate identically. Keep the macro, the residual plot, and
% every \boxguard byte-identical in the blank and the key.
\documentclass[12pt]{article}
\usepackage{apstats-article}
\usepackage{apstats-key}   % pulls in -boxes; do NOT also load -boxes
\newcommand{\answerspace}[2]{\par\nopagebreak\noindent\begin{minipage}[t][#1][t]{\linewidth}%
  \color{keyred}\bfseries #2\end{minipage}\par}

\begin{document}

\pageheader{Unit 2, Lesson 2.3}{Homework --- Answer Key}
% NAMESTRIP: no name row here — it lives on the cover only.

\vspace{0.08in}

% Authored BYTE-IDENTICALLY in homework/ and homework_key/.
\boxguard[8]
\begin{remindbox}
\small
\textbf{This is your graded homework} --- scored, and due the first class after two study halls.
\end{remindbox}

\vspace{0.02in}

\noindent A car dealer listed 18 used cars with 8 to 95 thousand miles. Price is in
thousands of dollars, mileage in thousands of miles:
\ $\widehat{\text{price}} = 32.0 - 0.24(\text{mileage})$. \ Answer every part \emph{in context}.

\vspace{0.02in}

\boxguard[14]
\begin{notesbox}{Part A --- Predicting with the Line}
\begin{enumerate}[leftmargin=1.9em, itemsep=2pt]
  \item Predict the price of a car with 50 thousand miles. Show your work.
        \answerspace{0.7cm}{$\widehat{\text{price}} = 32.0 - 0.24(50) = 20.0$, so about \$20{,}000.}
  \item A car with 50 thousand miles is listed at \$22{,}500. Find and interpret its residual.
        \answerspace{1.2cm}{Residual $= 22.5 - 20.0 = +2.5$: the car is listed \$2{,}500 above the line's prediction --- the line underpredicted its price.}
  \item Predict the price at 150 thousand miles. What went wrong, and why?
        \answerspace{1.2cm}{$32.0 - 0.24(150) = -4.0$: a negative price. 150 thousand miles is far outside the 8--95 in the data --- extrapolation.}
\end{enumerate}
\end{notesbox}

\vspace{0.02in}

\boxguard[16]
\begin{notesbox}{Part B --- Reading Residuals}
\begin{enumerate}[leftmargin=1.9em, itemsep=2pt]
  \item \begin{minipage}[t]{0.52\linewidth}
        Here is the residual plot for the 18 cars. Is a linear model appropriate for
        predicting price from mileage? Explain.
        \answerspace{1.2cm}{Yes: random scatter around 0, no pattern.}
        \end{minipage}\hfill
        \begin{minipage}[t]{0.44\linewidth}
        \vspace{-10pt}\centering
        \begin{tikzpicture}[xscale=0.058, yscale=0.28]
          \draw[->, thick] (0,-3.3) -- (0,3.4);
          \draw[thick] (0,-3.3) -- (102,-3.3);
          \draw[dashed, slate] (0,0) -- (102,0);
          \foreach \y in {-3,0,3} {\draw (-1.2,\y) -- (1.2,\y);
            \node[left, font=\tiny] at (0,\y) {\y};}
          \foreach \x in {0,20,40,60,80,100} {\draw (\x,-3.42) -- (\x,-3.18);
            \node[below, font=\tiny] at (\x,-3.3) {\x};}
          \foreach \x/\e in {8/1.2, 12/-0.8, 18/2.1, 22/-1.5, 27/0.4, 31/-2.2, 36/1.6, 40/-0.3,
                             45/0.9, 50/2.5, 54/-1.8, 60/-0.6, 64/1.1, 70/-2.4, 75/0.7,
                             81/-1.0, 88/1.9, 95/-1.8}
            {\fill[navy] (\x,\e) circle[x radius=1.1, y radius=0.22];}
          \node[font=\tiny] at (51,-5.0) {Mileage (thousands)};
          \node[rotate=90, font=\tiny] at (-10,0) {Residual (\$1000s)};
        \end{tikzpicture}
        \end{minipage}
  \item \textbf{AP-style.} A car's residual is $-1.8$. Which is \textbf{correct}? Circle one and
        justify it.

        \begin{tabularx}{\linewidth}{@{} X X @{}}
        \textbf{(A)} Listed \$1{,}800 above the prediction. &
        \textcolor{keyred}{\textbf{(B)} Listed \$1{,}800 below the prediction.} \\
        \textbf{(C)} The line predicted \$1{,}800 too low. &
        \textbf{(D)} It has 1{,}800 fewer miles than predicted. \\
        \textbf{(E)} A linear model does not fit this car. & \\
        \end{tabularx}
        \answerspace{1.0cm}{Residual $= y - \hat{y} = -1.8$ is negative, so the actual price is \$1{,}800 below the prediction.}
\end{enumerate}
\end{notesbox}

\vspace{0.02in}

% The next-lesson preview closes the packet, so it lives here rather than in AP Practice.
\boxguard[4]
\begin{spiralbox}
\small
\textbf{Next: Least-Squares Regression} --- where $a$ and $b$ come from (the line with the
smallest squared residuals), and what the slope, intercept, and $r^2$ mean in context.
\end{spiralbox}

\end{document}
